% Printed 458; principal PDF 465; Google PDF 469.
\par\clearpage
\noindent\makebox[0pt][l]{\fontsize{13}{16}\selectfont 458}\hfill{\small\textbf{HEIGHT MEASUREMENTS.}}\hfill\null\par\medskip
\begin{center}
{\large Table III.}\par\medskip
\begingroup\setlength{\tabcolsep}{12pt}\renewcommand{\arraystretch}{1.05}
\begin{tabular}{r|r@{\hspace{12pt}\vrule width .9pt\hspace{12pt}}r|r@{\hspace{12pt}\vrule width .9pt\hspace{12pt}}r|r}
&$+$&&$+$&&$+$\\
1,9&1&2,8&4&3,4&17\\
2,3&1&2,9&5&3,5&22\\
2,4&2&3,0&7&3,6&27\\
2,5&2&3,1&9&3,7&34\\
2,6&3&3,2&11&3,8&43\\
2,7&3&3,3&14&3,9&54\\
2,8&4&3,4&17&&\\
\end{tabular}\endgroup
\par\bigskip\bigskip
{\large Use of the Tables.}
\end{center}
\par\noindent\hspace*{\parindent}$t$, $t'$ temperature of the air; $T$, $T'$ temperature of the mercury (on the Réaumur scale); $b$, $b'$ barometer reading (in \textls[120]{any} unit of measure).

\par\noindent\hspace*{\parindent}Reduce $\log b$ and $\log b'$ by $10T$, $10T'$, respectively (regarded as units in the $5^{\mathrm{th}}$ decimal place), and subtract the logarithms thus corrected from one another; let the difference be $=u$. Add $\log u$ and $A$, after correcting the latter, if deemed necessary, by the second table (which, like the third, gives units in the $5^{\mathrm{th}}$ decimal place); let the sum be $=v$; this quantity receives a further small correction from Table III, of which $v$ itself is the argument. The $v$ thus corrected is the logarithm of the difference in height in metres. If the difference is wanted in toises, add another $9{,}71018$ to the logarithm.
\par\medskip
\noindent\hspace*{.5em}\begin{tabular}{l@{\quad}l@{\quad}l@{\quad}l@{\quad}l}
Example.&$t=15^\circ{,}3,$&$T=14{,}9,$&$b=735{,}581\,\mathrm{mm},$&altitude of the pole $=45^\circ$.\\
&$t'=3{,}2,$&$T'=7{,}8,$&$b'=537{,}203$&
\end{tabular}
\par\noindent\hspace*{1.5cm}\begingroup\setlength{\tabcolsep}{2pt}\renewcommand{\arraystretch}{1.2}
\begin{tabular}{r@{\ }c@{\ }r@{\quad}c@{\ }r@{\ }l}
$\log b$&$=$&$2{,}86663$&Corr.&$-149$&\\
$\log b'$&$=$&$2{,}73014$&»&$-78$&\\\cline{2-3}\cline{5-5}
&&$0{,}13649$&»&$-71$&\qquad .\\
$u$&$=$&$0{,}13578$&&&\\\cline{2-3}
$\log u$&$=$&$9{,}13284$&&&\\
$A$&$=$&$4{,}28408$&&&\\
Corr.&$=$&$0$&&&\\\cline{2-3}
$v$&$=$&$3{,}41692$&Corr.&$+18$&$.\ .\ .\ .\ 3{,}41710=\log2612{,}8\,\mathrm{m}.$
\end{tabular}\endgroup
\begin{center}\rule{3.5cm}{.4pt}\end{center}

% Unnumbered title [459]; principal PDF 466; Google PDF 470.
\clearpage\null\vspace*{6.7cm}
\begin{center}{\fontsize{24}{30}\selectfont\textls[160]{HELIOTROPE.}}\end{center}
\par\vspace*{11cm}\noindent\makebox[.86\linewidth][r]{58\textsuperscript{*}}\par
% Blank reverse [460]; principal PDF 467; Google PDF 471.
\clearpage\null
% Unnumbered opening [461]; principal PDF 468; Google PDF 472.
\clearpage\vspace*{2.7cm}
\begin{center}
\begin{tabular}{c}\hline
\small\textbf{Göttingische gelehrte Anzeigen.\quad No. 126.\quad 9 August 1821.\quad pp. 1249--1254.}\\\hline
\end{tabular}\par\bigskip
{\large [On the Heliotrope.]}
\end{center}
\par\noindent\hspace*{\parindent}Those familiar with higher geodesy know the difficulties of obtaining suitable targets for forming \textls[120]{large} triangles. In some regions there are no tall church towers in suitable positions, and the existing towers often do not offer the desired conditions for setting up the instruments and centring the measured angles. Nor are they always built as regularly as one would wish for the utmost precision. Specially built signal towers, apart from the expense in money and time of building them, resemble church towers in that they are difficult to see and sight on at considerable distances where they stand out against a nearby dark background, and if given a light colour, they show a variable and highly detrimental phase according to how the sun lights them. Indeed, even the most perfect signal towers, blackened and silhouetted against the sky, are not entirely free of a troublesome phase at \textls[120]{very} great distances when one sees a sunlit side and a side in shadow at the same time. Measurements at night with \textsc{Argand} lamps are indeed free of these errors, but have other inconveniences, particularly on mountains that are difficult to reach, which are too obvious to need discussion here.

\par\noindent\hspace*{\parindent}These considerations led Court Councillor \textsc{Gauss}, for the

% Printed 462; principal PDF 469; Google PDF 473.
\par\clearpage
\noindent\makebox[0pt][l]{\fontsize{13}{16}\selectfont 462}\hfill{\small\textbf{HELIOTROPE.}}\hfill\null\par\medskip
\noindent arc measurement to be carried out in the Kingdom of Hanover by royal command, to devise a new aid which could be used to advantage in large-scale triangulation, if only alongside the one just mentioned. Its success has far surpassed even his expectations.

\par\noindent\hspace*{\parindent}An investigation based on photometric considerations had already convinced him that sunlight reflected by even a very small plane mirror must still be strong enough, at the greatest distances that could occur in triangulation, to provide a superb target. To put this idea to use, an instrument had to be devised by means of which sunlight could be directed precisely in any required direction wherever needed. At the same time, it had to be possible to set up and operate the instrument easily anywhere, and the centre of the reflecting mirror had to remain absolutely stationary throughout the movements required to follow the moving sun, as it were. An instrument that \textls[120]{directs sunlight} in any desired direction, perfectly and in the way indicated, would seem most aptly to bear the name \textls[120]{heliotrope}, at least as aptly as two well-known products of the plant and mineral kingdoms.

\par\noindent\hspace*{\parindent}Further reflection on the subject led Court Councillor \textsc{Gauss} to two entirely different designs for a heliotrope. Our skilful Inspector \textsc{Rumpf} has already made an excellent instrument according to one design, and the same craftsman is now working on a second according to the other. A complete description, which would be unsuitable for these pages, will be given elsewhere. --- In order to make actual trials of the strength of reflected sunlight before the heliotrope just mentioned was finished, Court Councillor \textsc{Gauss} hit upon a third idea by which any reflecting sextant can be made into a fairly good heliotrope, or indeed used as a heliotrope with nothing added if it can be mounted on a good stand. Court Councillor \textsc{Gauss} had Inspector \textsc{Rumpf} attach a third plane mirror to the large mirror of a sextant so that its

% Printed 463; principal PDF 470; Google PDF 474.
\par\clearpage
\noindent\null\hfill{\small\textbf{HELIOTROPE.}}\hfill\makebox[0pt][r]{\fontsize{13}{16}\selectfont 463}\par\medskip
\noindent plane is perpendicular to the plane of the sextant and makes with the plane of the large mirror an angle equal to the complement of the angle between the line of sight (which, if necessary, must first be established by fitting fine crosshairs) and the plane of the small mirror. As soon as such a sextant is placed as if to measure the angular distance of an object from the centre (or indeed, if everything is well made and adjusted, from any point) of the sun's disc, whichever of the two is seen directly, the third mirror reflects the sunlight towards the object. If such a sextant is mounted on a good stand, someone with a little practice can readily send the reflected sunlight continuously to the desired point. If there is no third mirror, the large mirror itself can take its place by quickly advancing the alidade by the angle mentioned above (or nominally by twice that angle on the graduated arc) while keeping the plane completely fixed. This latter method can plainly produce only intermittent reflection; but with due care and a few small tricks that would take too long to describe here, a practised hand can keep the sunlight reflected towards the specified point for two minutes or more at a time. Both of these latter methods plainly have the small imperfection that, so long as the stand remains fixed, the centre of the reflecting mirror is not absolutely stationary. In most cases, however, this will hardly matter at all, and if necessary one could easily make small adjustments continuously or allow for it in the calculations. Its applicability is of course also restricted to the angular distance of the sun from the object permitted by the length of the graduated arc of the sextant.

\par\noindent\hspace*{\parindent}Before reporting the outcome of the trials made with the heliotrope, we note that all the mirrors used for reflection are 2 inches wide and $1\frac14$ inches high. Experience has confirmed what Court Councillor \textsc{Gauss} had already calculated on photometric grounds, namely that in even moderately favourable circumstances larger dimensions would be quite unnecessary, at least for geodetic use. With the second design mentioned above, however, one can
% Printed 464; principal PDF 471; Google PDF 475.
\par\clearpage
\noindent\makebox[0pt][l]{\fontsize{13}{16}\selectfont 464}\hfill{\small\textbf{HELIOTROPE.}}\hfill\null\par\medskip
\noindent also have a larger mirror fitted at will without otherwise increasing the dimensions of the instrument.
\begin{center}
\begin{tikzpicture}[x=.007cm,y=-.007cm,line cap=round,line join=round]
\path[use as bounding box] (0,0) rectangle (2030,910);
\draw[line width=.65pt] (655,75) circle[radius=.238cm];
\fill (655,75) circle[radius=.0301cm];
\draw[line width=.45pt] (702,125)--(1165,750);
\draw[line width=.45pt] (78,750)--(1165,750)--(1976,740);
\draw[line width=.65pt] (644,735)--(977,735)--(979,740)--(979,760)--(975,765)--(646,766)--(641,760)--(641,741)--cycle;
\draw[line width=.65pt] (608,743)--(641,741)--(642,759)--(610,759)--(606,754)--cycle;
\draw[line width=.5pt] (877,618)--(1413,858);
\draw[line width=.5pt] (1234,607)--(1102,868);
\draw[line width=.5pt,dash pattern=on 0pt off 2.7pt] (884,634)--(1403,867);
\draw[line width=.5pt,dash pattern=on 0pt off 2.7pt] (1239,618)--(1116,867);
\end{tikzpicture}
\end{center}
\par\noindent\hspace*{\parindent}At the distance from the local observatory to the Hohehagen, a principal triangulation point of the arc measurement (nearly 2 geographical miles), the light from both the heliotrope itself and the sextant adapted as a heliotrope was splendidly visible to the naked eye when the sun shone brightly. In the telescope of the theodolite it was really too strong, whereas the reflected light from a bright cloud alone provided the finest target imaginable. Of course, the reflected sunlight itself can easily be dimmed as desired by covering part of the mirror.

\par\noindent\hspace*{\parindent}At the distance from the Hils (another principal triangulation point) to the observatory's meridian mark, very nearly 5 geographical miles, the light from both instruments was likewise still splendidly visible to the naked eye, like a beautiful little star, and provided a superb target in the telescope of the theodolite. At times in misty air, when not a trace of either the ridge of the Hils or the signal tower built there could be seen through the telescope of the theodolite, the heliotrope light seemed to float like a magnificent star in the blue sky.

\par\noindent\hspace*{\parindent}The most important trials were made only in the last few days. Professor \textsc{Encke}, director of the Seeberg Observatory, had come here at the invitation of Court Councillor \textsc{Gauss} to learn how to use the reflecting sextant as a heliotrope without a third mirror, and to learn the

% Printed 465; principal PDF 472; Google PDF 476.
\par\clearpage
\noindent\null\hfill{\small\textbf{HELIOTROPE.}}\hfill\makebox[0pt][r]{\fontsize{13}{16}\selectfont 465}\par\medskip
\noindent little tricks required. He then went to the Inselsberg while Court Councillor \textsc{Gauss} began the measurements on the Hohehagen. The former sent sunlight at intervals to the Hohehagen with the sextant used as a heliotrope (distance 85000 metres or $11\frac12$ geographical miles), from which the sunlight was directed towards the Inselsberg with the heliotrope proper. The trials and observations continued from 19 to 29 July under alternately unfavourable and favourable conditions and had the best outcome one could wish for. The heliotrope light gave both observers the finest targets imaginable. It often appeared like a beautiful little star when the outline of the mountain could scarcely be seen, or could not be seen at all, through the same telescope. Sometimes one observer was in fog and rain while the heliotrope light from the other side shone through strongly. Indeed, on several occasions several of those present on the Hohehagen who had exceptionally sharp eyesight believed they could make out the point of light on the Inselsberg with the naked eye. We may add that the angular measurements themselves, which were made using the heliotrope light, gave a degree of agreement on both sides that could not have been expected from any other signal at so great a distance except under especially favourable circumstances.

\par\noindent\hspace*{\parindent}These experiences already establish beyond doubt that, when heliotrope light is used, there will be no limits to the size of the triangles that can be formed beyond those imposed by the curvature of the Earth.

\par\noindent\hspace*{\parindent}As the needs of higher geodesy prompted the development of this instrument, we confine ourselves here to recounting the above experiences, without pursuing for now the obvious prospect of a future and perhaps even more important use of a means capable of reaching so powerfully across distance for telegraphic signalling in war and peace.
\par\vfill\noindent\makebox[0pt][l]{IX.}\makebox[.86\linewidth][r]{59}\par

% Printed 466; principal PDF 473; Google PDF 477.
\clearpage
\noindent\makebox[0pt][l]{\fontsize{13}{16}\selectfont 466}\hfill{\small\textbf{HELIOTROPE.}}\hfill\null\par\bigskip\bigskip
\begin{center}
\begin{tabular}{c}\hline
\small\textbf{Astronomisches Jahrbuch für das Jahr 1825.\quad Berlin 1822.\quad pp. 103 and 104.}\\\hline
\end{tabular}\par\bigskip
{\large Invention of a Heliotrope,\\
by Court Councillor and Knight \textsc{Gauss} in Göttingen,\\
submitted under date of 26 December 1821.}
\end{center}
\par\noindent\hspace*{\parindent}You will excuse my long silence when I tell you that I was away for most of the year and even travelled to Altona late in the autumn to take delivery of \textsc{Ramsden}'s zenith sector, returning here only recently.

\par\noindent\hspace*{\parindent}In my triangulation, in which I have so far measured the angles at five triangulation points, I have sought to make the triangles as large as possible. You will have read the report on the new aid I have used for this purpose, the \textls[120]{heliotrope}, and on the first large-scale trials with it, in No. 126 of the local gelehrte Anzeigen. Since then I have used it continually, both as a target in measuring angles and, with no less success, for telegraphic signalling. The powerful effect of sunlight reflected from a mirror 2 inches wide and $1\frac14$ inches high, which was seen at distances of 5, 6, $7\frac12$, and even once $9\frac13$ geographical miles with the \textls[120]{naked eye}, usually astonishes those who experience it for the first time and have not been prepared for it by theoretical calculation. When the air is at all favourable, the size of the triangle sides is now limited only by the curvature of the Earth, especially if one gives the mirrors somewhat larger dimensions, as I have done with two newly made heliotropes of entirely different design.
\begin{center}\rule{3.5cm}{.4pt}\end{center}

% Printed 467; principal PDF 474; Google PDF 478.
\par\clearpage
\noindent\null\hfill{\small\textbf{CORRESPONDENCE.\quad HELIOTROPE.}}\hfill\makebox[0pt][r]{\fontsize{13}{16}\selectfont 467}\par
\vspace*{4.2cm}
\begin{center}{\fontsize{19}{24}\selectfont\textls[80]{CORRESPONDENCE.}}\end{center}
\vspace*{1.5cm}
\begin{center}{\large\textsc{Gauss} to \textsc{Olbers}.\quad Göttingen, 1 July 1821.}\end{center}
\par\noindent\hspace*{\parindent}.\ .\ .\ .\ .\ .\enspace Meanwhile, what mattered to me at first was simply to gain some experience of the strength of the light from the mirrors themselves. I first tried various things. I fastened a mirror to the cover of a theodolite telescope and sought, by means of azimuths and altitudes laboriously calculated in advance, to give the mirror the right position to cast the light in a particular direction. This failed completely, however. The cover, which moved rather stiffly, could not reliably be brought back to the same position every time. Differences of $20'$ remained, making this method quite unusable unless the mirror could be fastened more solidly to the telescope so that the latter stayed open. In the meantime, annoyance at the wasted effort led me to another idea, which succeeded completely. The reflecting sextant alone, on a good stand, does what is required, though not as well as a proper sun mirror. If the angle that the line of sight (which, if there is none already, must first be established by a thread or crosshairs) makes with the small mirror is $=90^\circ-a$, and if the image of the sun and the object towards which the light is to be cast are brought into contact in the usual way as if to measure the angle --- regardless of whether the former or the latter is seen directly ---, one need only advance the alidade by $a$ (or nominally by $2a$) while keeping the plane fixed to achieve the desired result. With a little practice one can readily set it so that the coincidence
\par\medskip\noindent\makebox[.86\linewidth][r]{59\textsuperscript{*}}\par
% Continuation on printed 468.
